\documentclass[12pt,a4paper,UTF8]{ctexart}
\usepackage[a4paper,top=22mm,bottom=22mm,left=22mm,right=22mm,headheight=15pt]{geometry}
\usepackage{amsmath,amssymb,mathtools,bm}
\usepackage{graphicx}
\usepackage{xcolor}
\usepackage{booktabs,longtable,array}
\usepackage{enumitem}
\usepackage{fancyhdr}
\usepackage{caption}
\usepackage{float}
\usepackage{microtype}
\usepackage{setspace}
\usepackage[hidelinks,unicode]{hyperref}
\usepackage[most]{tcolorbox}

\definecolor{ContestBlue}{HTML}{17365D}
\definecolor{SoftBlue}{HTML}{EEF4FA}
\definecolor{SoftGray}{HTML}{F5F6F7}
\hypersetup{pdftitle={第088套},pdfauthor={}}

\setstretch{1.24}
\setlength{\parindent}{2em}
\setlength{\parskip}{0.35em}
\setlist{itemsep=0.35em,topsep=0.4em}
\captionsetup{font=small,labelfont=bf,justification=centering}
\setcounter{secnumdepth}{0}

\pagestyle{fancy}
\fancyhf{}
\fancyhead[L]{\small 第088套}
\fancyfoot[C]{\thepage}
\renewcommand{\headrulewidth}{0.4pt}

\ctexset{
  section={format=\Large\bfseries\color{ContestBlue},beforeskip=1.2em,afterskip=0.7em},
  subsection={format=\large\bfseries\color{ContestBlue},beforeskip=1.0em,afterskip=0.5em},
  subsubsection={format=\normalsize\bfseries,beforeskip=0.8em,afterskip=0.35em}
}

\makeatletter
\newcommand{\tightlist}{\setlength{\itemsep}{0pt}\setlength{\parskip}{0pt}}
\newcommand{\maxwidth}{\ifdim\Gin@nat@width>\linewidth\linewidth\else\Gin@nat@width\fi}
\newcommand{\maxheight}{\ifdim\Gin@nat@height>0.72\textheight 0.72\textheight\else\Gin@nat@height\fi}
\newcommand{\pandocbounded}[1]{\begingroup\setkeys{Gin}{width=\maxwidth,height=\maxheight,keepaspectratio}#1\endgroup}
\makeatother

\begin{document}

\begin{titlepage}
  \centering
  \vspace*{58mm}
  {\Huge\bfseries\color{ContestBlue} 第088套\par}
  \vfill
  \vspace*{28mm}
\end{titlepage}


\section{第 1 题}\label{ux7b2c-1-ux9898}

如图所示，一个质量为 \(m\)、半径为 \(r\)
的匀质薄壁圆筒由轴承限制，只能绕通过圆心 \(O\)
的中心轴转动。轴承铅垂安装在水平面上，圆筒质心位于水平光滑桌面内。一根长度为
\(l=2r\) 的不可伸长细绳，一端系在圆筒边缘的 \(B\) 点，另一端连接质量为
\(m\) 的质点 \(A\)。系统初始静止。沿垂直于细绳的方向对质点施加大小为

\[
J=mv_0
\]

的冲量。不计重力。

试判断细绳与系绳点处法线 \(OB\) 的夹角能否达到
\(30^\circ\)。若能，求此时圆筒相对地面的角速度；若不能，说明理由。

\begin{figure}[H]
\centering
\includegraphics[width=0.9\linewidth,height=\textheight,keepaspectratio,alt={图 1-1　圆筒、细绳和质点的初始构型。}]{../figures/question-01.pdf}
\caption*{图 1-1　圆筒、细绳和质点的初始构型。}
\end{figure}

\newpage

\section{第 2 题}\label{ux7b2c-2-ux9898}

万有引力常量为 \(G\)。三个质量均为 \(m\) 的星体保持边长为 \(l\)
的正三角形相对构型，在万有引力作用下绕质心以恒定角速度 \(\omega\) 转动。

\begin{enumerate}
\def\labelenumi{\arabic{enumi}.}
\tightlist
\item
  求 \(\omega\)，以及体系总相互作用势能 \(E_p\) 与总动能 \(E_k\) 的比值
  \(E_p/E_k\)；
\item
  将三个星体的质量改为 \(m,\,2m,\,3m\)，重新求 \(\omega\) 与
  \(E_p/E_k\)；
\item
  将体系改为四个质量均为 \(m\) 的星体，保持边长为 \(l\)
  的正方形相对构型，求 \(\omega\) 与 \(E_p/E_k\)；
\item
  考虑由 \(N\) 个质量均为 \(m\) 的星体组成的正 \(N\)
  边形，相邻星体间距为 \(l\)。求 \(\omega\) 与 \(E_p/E_k\)。
\end{enumerate}

\begin{figure}[H]
\centering
\includegraphics[width=0.9\linewidth,height=\textheight,keepaspectratio,alt={图 2-1　三角形与正方形相对平衡构型。}]{../figures/question-02.pdf}
\caption*{图 2-1　三角形与正方形相对平衡构型。}
\end{figure}

\newpage

\section{第 3 题}\label{ux7b2c-3-ux9898}

如图所示，两个绝热圆筒的底面积分别为 \(2S\) 和
\(S\)，底部通过一个细的绝热阀门连通。左右圆筒内各有一个密闭、绝热、不计阻力和重力的活塞，活塞上端通过理想轻绳和定滑轮连接。两活塞上方均为大气，压强为
\(P_0\)。初始时阀门关闭，左缸充有体积为 \(2V_0\)、温度为 \(T_0\)
的理想气体，右缸为体积 \(V_0\) 的真空。气体的定容摩尔热容为

\[
C_V=\frac32R.
\]

\begin{enumerate}
\def\labelenumi{\arabic{enumi}.}
\tightlist
\item
  缓慢打开阀门，研究左右气体达到压强平衡的过程，求相应的左右体积
  \(V_L,V_R\) 和温度 \(T_L,T_R\)；
\item
  将活塞改为导热活塞，并将左缸截面积改为 \(xS\)、左缸初始体积改为
  \(xV_0\)。打开阀门后，平衡状态下左缸体积是否可能不为零？若可能，求
  \(x\) 的范围。
\end{enumerate}

提示：活塞导热率很低，在机械平衡已建立但热平衡尚未建立时，两侧气体可在较长时间内维持不同温度。理想气体绝热方程为

\[
pV^\gamma=C,
\qquad
\gamma=\frac{C_V+R}{C_V}.
\]

\begin{figure}[H]
\centering
\includegraphics[width=0.611\linewidth,height=\textheight,keepaspectratio,alt={图 3-1　双活塞圆筒与底部阀门。}]{../figures/question-03.pdf}
\caption*{图 3-1　双活塞圆筒与底部阀门。}
\end{figure}

\newpage

\section{第 4 题}\label{ux7b2c-4-ux9898}

两个静止质量均为 \(m_0\) 的探测器，在地面参考系中以速度 \(v_0=0.8c\)
相向运动。为避免相撞，其中一个探测器沿前进方向发射大量光子，这些光子被另一个探测器完全吸收，最终两个探测器相对静止。

\begin{enumerate}
\def\labelenumi{\arabic{enumi}.}
\tightlist
\item
  求地面系中所发射光子的总能量 \(E_0\)；
\item
  当地面系中累计发射的光子能量为 \(E\) 时，求发射器此时的静止质量 \(m\)
  和地面系速度 \(v\)；
\item
  若在发射器的瞬时共动参考系中，单位时间发射光子的功率恒为
  \(W\)，求地面系中完成全部发射所需的时间 \(t\)。
\end{enumerate}

提示：

\[
\int_a^b\frac{dx}{x}=\ln\frac ba.
\]

\begin{figure}[H]
\centering
\includegraphics[width=0.9\linewidth,height=\textheight,keepaspectratio,alt={图 4-1　相向运动的探测器与光子传能。}]{../figures/question-04.pdf}
\caption*{图 4-1　相向运动的探测器与光子传能。}
\end{figure}

\newpage

\section{第 5 题}\label{ux7b2c-5-ux9898}

如图所示，一个正八面体框架的每条棱长均为 \(l\)，每条棱的电阻均为
\(r\)。电流 \(I\) 从顶点 \(A\) 流入，从与 \(A\) 相邻的顶点 \(B\) 流出。

\begin{enumerate}
\def\labelenumi{\arabic{enumi}.}
\tightlist
\item
  求 \(A,B\) 两点间的电势差；
\item
  空间中存在沿正八面体上下对角线方向的匀强磁场
  \(\bm B\)，求八面体所受支撑力的合力。
\end{enumerate}

\begin{figure}[H]
\centering
\includegraphics[width=0.462\linewidth,height=\textheight,keepaspectratio,alt={图 5-1　正八面体电阻网络、电流端点与磁场方向。}]{../figures/question-05.pdf}
\caption*{图 5-1　正八面体电阻网络、电流端点与磁场方向。}
\end{figure}

\newpage

\section{第 6 题}\label{ux7b2c-6-ux9898}

一个静止质量为 \(m\)、电量为 \(q\)
的电荷初始静止在原点。空间存在匀强电场和匀强磁场

\[
\bm E=(0,-E_0,0),
\qquad
\bm B=(0,0,-B_0),
\]

且满足

\[
E_0=cB_0,
\]

其中 \(c\) 为光速。考虑狭义相对论效应，求粒子的运动轨迹，并用
\(x,y,m,c,q,E_0\) 表示。

\begin{figure}[H]
\centering
\includegraphics[width=0.86\linewidth,height=\textheight,keepaspectratio,alt={图 6-1　相互垂直的匀强电场和磁场。}]{../figures/question-06.pdf}
\caption*{图 6-1　相互垂直的匀强电场和磁场。}
\end{figure}

\newpage

\section{第 7 题}\label{ux7b2c-7-ux9898}

如图所示，一条均匀链条的长度为

\[
L=\frac{\pi R}{3},
\]

单位长度质量为 \(\lambda\)。链条贴在半径为 \(R\)
的光滑球面外侧，顶端固定在球面最高点。现从静止开始释放链条顶端。

\begin{enumerate}
\def\labelenumi{\arabic{enumi}.}
\tightlist
\item
  刚释放时，链条上各点的切向加速度大小是多少？
\item
  刚释放时，链条上何处的张力最大？
\item
  当链条顶端沿球面转过 \(\varphi=10^\circ\)
  且链条仍未离开球面时，求最大张力及其所在位置，并求该处单位长度链条对球面的压力。
\end{enumerate}

\begin{figure}[H]
\centering
\includegraphics[width=0.9\linewidth,height=\textheight,keepaspectratio,alt={图 7-1　光滑球面上的均匀链条。}]{../figures/question-07.pdf}
\caption*{图 7-1　光滑球面上的均匀链条。}
\end{figure}

\newpage

\section{第 8 题}\label{ux7b2c-8-ux9898}

容器中装有冰水混合物。在 \(t=0\ \mathrm{min}\) 时接通功率为

\[
P_0=400\ \mathrm W
\]

的电热炉。图中给出了物质温度 \(T\) 与时间 \(t\) 的关系。散热功率 \(Q\)
与温差 \(\Delta T=T-T_0\) 成正比，环境温度取
\(T_0=0^\circ\mathrm C\)，故

\[
Q=\alpha T,
\]

其中 \(\alpha\) 为与温度无关的常数。根据图像求：

\begin{enumerate}
\def\labelenumi{\arabic{enumi}.}
\tightlist
\item
  初始冰水混合物中冰的质量 \(m_{\rm ice}\)；
\item
  容器中物质的总质量 \(m\)；
\item
  散热比例系数 \(\alpha\)；
\item
  使水始终不会沸腾的最大加热功率 \(P_{\max}\)；
\item
  若改用功率 \(P_1=300\ \mathrm W\)
  的电热炉，求从冰开始熔化到水达到沸点所需的时间。
\end{enumerate}

已知水的比热容

\[
c=4200\ \mathrm{J/(kg\cdot{}^\circ C)},
\]

冰的熔解热

\[
\lambda=3.2\times10^5\ \mathrm{J/kg}.
\]

\begin{figure}[H]
\centering
\includegraphics[width=0.6\linewidth,height=\textheight,keepaspectratio,alt={图 8-1　冰水混合物的温度---时间曲线。}]{../figures/question-08.pdf}
\caption*{图 8-1　冰水混合物的温度---时间曲线。}
\end{figure}

\end{document}
